\documentclass[11pt]{amsart}
\usepackage[margin=1in]{geometry}
\usepackage{amssymb,mathtools,booktabs,longtable,url}
\usepackage[colorlinks=true,linkcolor=blue,citecolor=blue,urlcolor=blue,hypertexnames=false]{hyperref}
\newtheorem{theorem}{Theorem}
\newtheorem{proposition}[theorem]{Proposition}
\newtheorem{lemma}[theorem]{Lemma}
\newtheorem{corollary}[theorem]{Corollary}
\theoremstyle{remark}
\newtheorem{remark}[theorem]{Remark}
\newcommand{\degen}{\mathrel{\rightsquigarrow}}
\DeclareMathOperator{\rank}{R}

\title[The axioms behind $\omega\le 9/4$ stop at $9/4$]{Remark on OpenAI's $\omega\le 9/4$:\\
the profile axioms stop exactly at $9/4$, and an enlarged constraint system still does}
\author{[AUTHOR NAME]}
\date{October 9, 2026}

\begin{document}

\begin{abstract}
OpenAI's preprint ``An Upper Bound of $9/4$ for the Matrix Multiplication
Exponent'' (openai/math, family~107) obtains $\omega\le 9/4$ from six
properties (positivity, symmetry, boundary values, concavity, tripling, and
a rank bound) of a symmetrized profile $P(a,b)$ of the polynomial-multiplication
tensors. We give a self-contained proof that these six properties cannot
yield any exponent below $9/4$. The profile
$P^*(a,b)=\phi(m)\bigl(M+\tfrac{m-1}{2}\bigr)$, with $m=\min(a,b)$,
$M=\max(a,b)$ and $\phi(m)=\prod_{j<m}(1+\tfrac1{3j})$, is the pointwise least
profile satisfying the first five properties. It satisfies the rank
bound exactly at $t=3/4$. The least-element observation itself is not new: it
was recorded publicly on 2026-10-07 in the GitHub repository
\texttt{eliottcassidy2000/math} (THM-4568). It was also noted numerically on
2026-10-08 in \texttt{GasStationManager/SafeVerifyAgent}. Our proof of the
rank bound uses only the elementary inequality $\phi(m)^3\le 2m-1$.

The second part reports joint work with an OpenAI Codex session. A
\emph{balanced assignment}
$\lambda^*(C(a,b)^\pi)=P^*(a,b)^{3/4}$, $\lambda^*(\text{dot}_m)=m^{3/4}$,
$\lambda^*(T_m)=m^{9/4}$ satisfies, at all sizes, a substantially larger
system of valid character inequalities. These include several new
degenerations of polynomial multiplication. Hence that system cannot force
$\omega<9/4$ either. We state precisely which families are proved for all
sizes, which are checked only on finite ranges, and which routes remain open.
$\lambda^*$ is \emph{not} claimed to be a tensor character.
\end{abstract}

\maketitle

\section{Setting and attribution}

All tensor constructions and the proof of $\omega\le 9/4$ are OpenAI's
work~\cite{OAI}. We quote line numbers of \texttt{build/paper.tex} at
openai/math commit \texttt{adc7f1241b42}, which is unchanged at
\texttt{08fcd618213a}. For a character $\lambda$ with mean dot-product
exponent $t=(p_X+p_Y+p_Z)/3$, the paper defines the symmetrized profile of
$C(a,b)=\sum_{i<a,j<b}x_iy_jz_{i+j}$ (l.~420--422):
\[
P(a,b)=\Bigl(\prod_{\pi\in S_3}\lambda_\pi(C(a,b))\Bigr)^{1/(6t)}.
\]
The passage from characters to the exponent uses exactly the following
properties of $P$ and nothing else.
\begin{itemize}
\item[(Pos)] $P(a,b)>0$ (l.~427, l.~611).
\item[(S)] $P(a,b)=P(b,a)$ (l.~427, l.~611).
\item[(B)] $P(1,b)=b$ (l.~427, l.~611).
\item[(C)] $2P(a,b)\ge P(a,b+1)+P(a,b-1)$ for $a\ge1,b\ge2$ (Lemma~4.1,
l.~442--447; hypothesis at l.~612).
\item[(T)] $P(a,3h+a-1)\ge 3P(a,h)$ for $a,h\ge1$ (Lemma~4.2, l.~526--530;
hypothesis at l.~613).
\item[(U$_t$)] $P(a,b)\le(a+b-1)^{1/t}$ (l.~428). Only the diagonal case
$P(a,a)\le(2a-1)^{1/t}$ is used (l.~657).
\end{itemize}
Lemma~5.1 (Diagonal growth, l.~610--650) takes (Pos), (S), (B), (C), (T) as
its only hypotheses. The proof of the main theorem (l.~652--660) adds only
(U$_t$). The character-existence lemma (l.~198--202, Appendix~A) is used
afterwards, only to pass from ``$t\le 3/4$ for every character'' to
$\nu\le 9/4$. It places no further constraint on~$P$.

The Lean formalization in the same repository has the same structure.
\path{structure ScalarProfile} has exactly the six fields above
(\path{Growth/Profile.lean}, l.~23--33). The theorem
\path{meanExponent_le_three_quarters} reaches the exponent only through
\path{toScalarProfile} and \path{ScalarProfile.exponent_le_three_quarters}
(\path{Polynomial/Inequalities.lean}, l.~97--114). These paths are relative
to \path{lean/OAI/LinearAlgebra/MatrixMultiplication/AuxiliarySeparation/}.
\emph{Related work.} Independent formal checks of~\cite{OAI} include
\cite{SEL}, which extends the proof to all fields, and \cite{SVA}. The
latter also observes numerically (\S2d) that the least feasible diagonal
equals the bound in the paper's proof. Earlier ``limits of the method''
results concern the laser method and Coppersmith--Winograd tensors
\cite{AFLG,Alm,CVZb}. They do not apply to the profile argument, which is
built on the spectral framework \cite{Str88,CVZ,LG14}. This remark
complements \cite{SEL}: \cite{SEL} shows that the $9/4$ argument works over
every field, and we show that its profile axioms cannot go below $9/4$.

We write $\mathrm{Ax}(t)$ for the system (Pos), (S), (B), (C), (T), (U$_t$).
Every character satisfies $1/3\le t\le 1$ (l.~211--234).

\section{The least profile}

Put $\phi(1)=1$, $\phi(m)=\prod_{j=1}^{m-1}\bigl(1+\frac1{3j}\bigr)
=\frac{\Gamma(m+1/3)}{\Gamma(4/3)\Gamma(m)}$, and
\[
P^*(a,b)=\phi(m)\Bigl(M+\frac{m-1}{2}\Bigr),\qquad m=\min(a,b),\ M=\max(a,b).
\]
Its diagonal values are $1,\frac{10}3,\frac{56}9,\frac{770}{81},\dots$, and
$P^*(a,a)=\frac{3a-1}{2}\phi(a)\sim\frac{3}{2\Gamma(4/3)}a^{4/3}$.

\begin{theorem}[least profile; recorded earlier in \cite{EC}]\label{thm:main}
The following hold.
\begin{enumerate}
\item[(i)] $P^*$ satisfies (Pos), (S), (B), (C), (T), and also (U$_{3/4}$).
Equality holds in (U$_{3/4}$) only at $(1,1)$.
\item[(ii)] Every $P$ satisfying (Pos), (S), (B), (C), (T) satisfies
$P\ge P^*$ pointwise.
\item[(iii)] Hence, for $t>0$, $\mathrm{Ax}(t)$ is satisfiable if and only if $t\le 3/4$.
In particular no argument that uses only these properties of the profile can
give $t<3/4$, that is, $\omega<9/4$. Every inequality in the proof of
Lemma~5.1, up to and including \eqref{eq:H}, is an equality for $P^*$.
\end{enumerate}
\end{theorem}

\begin{proof}
(i) (Pos), (S) and (B) are immediate.

\emph{(C).} Fix $a$ and put $f(b)=P^*(a,b)$. For $b\ge a$, $f$ is linear
with slope $\phi(a)$. For $b+1\le a$, the increment is
$\delta(b)=f(b+1)-f(b)=\phi(b)\bigl(\frac{a}{3b}+\frac23\bigr)$. The
inequality $\delta(b+1)\le\delta(b)$ is equivalent to
$(3b+1)(a+2b+2)\le(a+2b)(3b+3)$, which is $2(a-b-1)\ge0$. This holds for
$b\le a-2$, which is the range needed. At the junction,
$\delta(a-1)=\phi(a-1)\frac{3a-2}{3(a-1)}=\phi(a)$, so the increments are
non-increasing everywhere. They are also positive.

\emph{(T).} Since $3h+a-1\ge a$, the left side is
$\phi(a)\bigl(3h+\frac32(a-1)\bigr)$. If $h\ge a$, the right side is
$3\phi(a)(h+\frac{a-1}2)$, so equality holds. If $h<a$, put
$L(a')=\phi(a')(3h+\frac32(a'-1))$ and $R(a')=3\phi(h)(a'+\frac{h-1}{2})$ for
$a'\ge h$. Then $L(h)=R(h)$, and
\[
L(a'+1)-L(a')=\phi(a')\Bigl(2+\frac{h}{a'}\Bigr),\qquad R(a'+1)-R(a')=3\phi(h).
\]
At $a'=h$ the two increments are equal. The left increment is non-decreasing
in $a'$: its forward difference is
$\phi(a')\frac{2(a'+1-h)}{3a'(a'+1)}\ge0$. Hence $L(a)\ge R(a)$, with equality
exactly when $h\ge a-1$.

\emph{(U$_{3/4}$).} By induction, $\phi(m)^3\le 2m-1$. The base case is
$m=1$. The inductive step $(1+\frac1{3m})^3\le\frac{2m+1}{2m-1}$ is
equivalent to $(2m+1)(3m)^3-(2m-1)(3m+1)^3=9m^2+7m+1>0$. Since
$M+m-1\ge M+\frac{m-1}2$ and $M+m-1\ge 2m-1$,
\[
P^*(a,b)\le\Bigl(M+\frac{m-1}{2}\Bigr)(2m-1)^{1/3}\le (a+b-1)^{4/3}.
\]
For $m\ge2$ the first inequality is strict. For $m=1$ and $M\ge2$ the second
is strict. Alternatively, Wendel's inequality
$\Gamma(m+\frac13)/\Gamma(m)\le m^{1/3}$~\cite{Wendel} gives
$\phi(m)^3\le\Gamma(4/3)^{-3}m\approx1.4043\,m\le 2m-1$ for $m\ge2$.

(ii) We follow the paper's own chain (l.~617--639). For fixed $a$, (C) makes
$\Delta_{a,h}=P(a,h+1)-P(a,h)$ non-increasing in $h$. Write $D_a=P(a,a)$.
Lemma~5.1's proof gives
\begin{equation}\label{eq:H}
D_a\ \ge\ \frac{3a-1}{2}\prod_{m=1}^{a-1}\Bigl(1+\frac1{3m}\Bigr)=P^*(a,a).
\end{equation}
Off the diagonal, take $b>a$ and $h\ge a$. By (T) and (C),
\[
2P(a,h)\le P(a,3h+a-1)-P(a,h)=\sum_{k=h}^{3h+a-2}\Delta_{a,k}\le(2h+a-1)\Delta_{a,h},
\]
so $P(a,h+1)\ge P(a,h)\frac{2h+a+1}{2h+a-1}$. Telescoping from $h=a$ to
$b-1$ gives
\[
P(a,b)\ge D_a\cdot\frac{2b+a-1}{3a-1}\ge P^*(a,b).
\]
By (S) the same bound holds for $b<a$. This off-diagonal step was also given
by an auditor in~\cite{EC} and in~\cite[Codex round~3, (1)]{SUPP}.

(iii) If $0<t\le3/4$, then $(a+b-1)^{1/t}\ge(a+b-1)^{4/3}$, so $P^*$
satisfies $\mathrm{Ax}(t)$ by~(i). Conversely, Lemma~5.1 gives
$a^{4/3}\le P(a,a)\le(2a-1)^{1/t}$ for every $a$, which forces $t\le 3/4$.
The tightness statements follow from the equality cases in~(i): $P^*$
satisfies $\Delta_{a,a}=\phi(a)$, equality in (T) at $h=a$, and equality in
the diagonal recursion. Only the last two simplifications in the paper,
$(1+\frac1{3m})^3\ge1+\frac1m$ and $\frac{3a-1}{2}\ge a$, are strict. They
lose the constant $\frac3{2\Gamma(4/3)}\approx1.68$ but not the exponent.
\end{proof}

\section{What the barrier does and does not say}

\begin{enumerate}
\item It concerns the \emph{axiom system} $\mathrm{Ax}(t)$, not tensors.
Theorem~\ref{thm:main} does not say that some character has profile $P^*$.
It does not say that $\omega=9/4$. It does not say that polynomial
multiplication tensors can never prove $\omega<9/4$.
\item To go below $9/4$ with this approach, one needs a further valid
inequality for characters that every profile in the feasible set at
$t=3/4$ violates. That set contains $P^*$ and, by (ii), only profiles
$\ge P^*$. Alternatively, one needs information that the symmetrized profile
does not see.
\item The formal values below are not claimed to come from a character. If
some character $\lambda$ had $\lambda(T_m)=m^{9/4}$ for some $m\ge2$, then by
multiplicativity and $\lambda\le\rank$ we would get
$\rank(T_{m^k})\ge m^{9k/4}$ for all $k$, so $\omega\ge 9/4$. Together with~\cite{OAI} this
would give $\omega=9/4$; compare Strassen's spectral characterisation of
$\omega$~\cite{Str88}. Nothing here proves or refutes this.
\end{enumerate}

\section{An enlarged constraint system (joint work with OpenAI Codex)}\label{sec:enlarged}

In three rounds on 2026-10-09, an OpenAI Codex session (model
\texttt{gpt-6-astra}, Codex CLI v0.157.0) and we searched for \emph{valid}
character inequalities that might exclude the balanced assignment
\begin{gather*}
p_X=p_Y=p_Z=t=\tfrac34,\qquad
\lambda^*(C(a,b)^\pi)=P^*(a,b)^{3/4}\ (\pi\in S_3),\\
\lambda^*(B_L(m))=m^{3/4},\qquad \lambda^*(\langle e,h,l\rangle)=(ehl)^{3/4}.
\end{gather*}
The constructions in rows A--D, F, the two scalar-target rows, the
single-block row, the $T_m$ rows, CW and zero/rest of
Table~\ref{tab} were proposed and proved by Codex. The scope discipline in
the table is also Codex's. We supplied $P^*$ as the witness; it had been found
earlier and independently in~\cite{EC}. We also supplied the unequal-$h$
chain and Kronecker rows; the equal-$h$ chain and its tangent-line argument
are already in~\cite{EC} (``gadget barrier''). We independently re-checked
every row as indicated. The transcripts are
reproduced in~\cite{SUPP}.

\begin{theorem}[enlarged barrier; Codex and us]\label{thm:enlarged}
$\lambda^*$ satisfies every inequality of every family marked \textbf{P} in
Table~\ref{tab}, for all sizes and parameters. It also satisfies every
consequence obtained by adding or multiplying these numerical inequalities,
by optimizing the entropy vector $q$ separately for each leg permutation, or
by mixing them in a linear program. Hence this system cannot force $t<3/4$.
\end{theorem}

Here ``valid'' means that the inequality holds for \emph{every} character,
via an explicit restriction or degeneration together with additivity,
multiplicativity and the shared-leg inequality (paper Cor.~3.2). The
validity proofs are in the paper, in~\cite{SUPP}, or (for odd-$k$
chains) the following weighting. Give the even-position sectors, which are
copies of $C$, weight $0$. Give the odd-position sectors, which are copies of
$C^\sigma$, weight $+1$ on $Y$ and $-1$ on $Z$. This generalizes l.~544--548.
The ``tensor'' column records only finite coefficient or support checks,
which add evidence but are not proofs. The status column concerns ``$\lambda^*$
satisfies the family''. An invalid inequality that $\lambda^*$ satisfies
would not weaken the barrier. The
balanced lift turns each nonsymmetrized entropy inequality into its
symmetrized form: with all $\lambda_\pi$ equal, optimizing
$e^{\frac34H(q)}\prod_iF_i^{\frac34q_i}$ gives $(\sum_iF_i)^{3/4}$.

{\small
\begin{longtable}{p{0.08\textwidth}p{0.46\textwidth}p{0.06\textwidth}p{0.29\textwidth}}
\caption{Valid inequality families and the status of ``$\lambda^*$
satisfies them''. \textbf{P}: proved for all sizes. \textbf{E}: exact check
on a finite range only. \textbf{IC}: interval-certified finite recipe
search. Tensor = verification of the underlying degeneration.}\label{tab}\\
\toprule
id & inequality (source tensor relation) & status & evidence \\ \midrule
paper & (C), (T) and their unsymmetrized entropy forms, any $q$ per $\pi$ &
\textbf{P} & Thm.~\ref{thm:main}; balanced-lift reduction\\
chain & $P(a,\sum_{i\le k}h_i+\frac{(k-1)(a-1)}2)\ge\sum P(a,h_i)$ (odd
$k$, alternating $C$/$C^\sigma$ sectors sharing $X$; $k=3$, equal $h_i$ is
Lemma~4.2) & \textbf{P} & tangent line $\phi(a)(b+\frac{a-1}2)\ge P^*(a,b)$
(equal $h_i$: \cite{EC}); E for $k=3,5$; tensor ($k\le7$): ours\\
A & $mP(a,b)\ge\sum_{i\in I}P(a+m-1-2i,b-m+1+2i)$ (via
$C(a,b)\otimes B_Z(m)$, apolar filtration) & \textbf{P} & Jensen along
$a+b=$const; slack formula re-checked (E); tensor: Codex (780 coeff.)\\
B & $mP(a,b)\ge\sum_{i<\min(b,m)}P(a,b+m-1-2i)$ (via $C(a,b)\otimes B_X(m)$;
$m=2$ is Lemma~4.1) & \textbf{P} & Jensen; tensor: ours + Codex\\
C & $P(a,b)P(c,d)\ge\sum_{i+j=k}P(a+c-1-2i,b+d-1-2j)$ (bivariate
$D$-filtration, shared output) & \textbf{P} & proof in Codex round~2; E;
tensor: ours + Codex\\
D & $[P(a,b)P(c,d)]^t\ge\sum_{i<r}P(a+c-1-2i,b+d-1-2i)^t$, $r=\min(a,b,c,d)$ (full direct sum)
& \textbf{P} & rational certificate $[\,\cdot\,]^3\ge rS^3$; E; tensor: ours\\
F & $(2M-1)P(Ma,Mb)\ge MP(a,b)$ (Fourier filtering, cost kept) & \textbf{P}
& monotonicity\\
O,L & $\max(a,b)\le P(a,b)\le ab$ (outer product; dots) & \textbf{P} &
$\phi(a)\le 2a/3$\\
wrap & $P(n,n)^t\ge n$ ($C(n,n)\ge\mathbb C[\mathbb Z_n]\cong\langle
n\rangle$) & \textbf{P} & $\phi(n)^3\ge n$\\
peel & $P(n,n)^t\ge1+P(n-1,n-1)^t$ ($C(n,n)\degen1\oplus C(n-1,n-1)$) &
\textbf{P} & $f_n/n$ non-decreasing; tensor: ours\\
CW & $P(n,n)^t\ge(27/4)^{1/3}(n-2)^t$ ($C(n,n)\degen CW_{n-2}\ge cw_{n-2}$,
standard laser) & \textbf{P} & AM--GM; tensor (support): ours\\
scalar & any tensor product of generators $\degen\langle N\rangle$: $N\le
\lambda^*$ & \textbf{P} & $(d_Xd_Yd_Z)^{1/3}\le P^{*3/4}$\\
1-block & monomial degeneration (in convolution/group coordinates) of any
product of $C$'s, dots and matrix factors to \emph{one}
$\langle e,h,l\rangle$ & \textbf{P} & abelian-group lemma $ehl\le|G|$\\
$T\to C$ & $T_m\degen\bigoplus_sC(a_s,b_s)^{\pi_s}$ (any basis change) &
\textbf{P} & flattening sandwich, $m^{1/4}>9/8$\\
$T\otimes C$ & $T_m\otimes C(n,n)\degen\bigoplus C$'s, $m\ge2$ &
\textbf{P} & sufficient condition (24) of round~2\\
0/rest & laser on $\{000,011,101,112\}$ partition of $C(n,n)$, any child
value $\le$ border rank & \textbf{P} & entropy dual; head $q\le15$ interval,
tail rational\\
Kron & $P(a+(a'-1)N,b+(b'-1)N)\ge P(a,b)P(a',b')$, $N=a+b-1$ & E & $a,b,a',b'\le9$\\
interval & recursive consecutive-interval laser recipes, $n\le30$; all root
cuts, $n\le12$ & IC & Codex round~3; re-run by us\\
\bottomrule
\end{longtable}
}

\begin{remark}[open routes]
The following are \emph{not} covered by Theorem~\ref{thm:enlarged}:
\begin{itemize}
\item full direct sums of several matrix blocks extracted from powers of
$C(n,n)$, especially after non-monomial changes of basis;
\item higher-power repartitionings other than the families above;
\item $T\otimes C$ relations outside the sufficient condition (for example
pure convolution, or very unbalanced $(a,b)$), and targets that contain
products of several $C$'s or matrix factors;
\item relations that genuinely use asymmetry ($\lambda_\pi$ unequal) at
$t=3/4$.
\end{itemize}
The uniform quantum functionals give
$\mathcal Q(C(a,b))=(ab(a+b-1))^{1/3}$, but on matrix tensors they have total
degree $2$ rather than $9/4$. So they do not close the first route
(Codex round~3, \S7).
\end{remark}

\section{Computational certificates}
The repository's \texttt{scripts/} directory contains four programs, each
with its expected output. All decisions use exact rationals, except where a
floating-point step is noted.
\begin{itemize}
\item \texttt{r2\_profile\_axioms.py} checks (i) exhaustively for
$a,b\le200$ ($b<800$ for (C)). It also checks the equality cases, the
tightness of every step of Lemma~5.1, both induction identities, and Wendel's
bound with an interval-certified constant.
\item \texttt{r2\_extension\_families.py} is an independent re-implementation
of the inequality checks for Table~\ref{tab}, about $0.9\times10^6$ exact
comparisons in total. The $3/4$-powers are bracketed by integer fourth roots.
\item \texttt{r2\_tensor\_constructions.py} verifies the underlying
degenerations at the coefficient level for small sizes. These are the paper's
tripling degeneration, the alternating $k$-sector chains ($k\le7$), peeling, the CW and zero/rest supports, and the
$D$-adapted bivariate products ($7600$ basis-pair products) behind B, C, D
and Lemma~4.1. A mutation test confirms that the checks detect deliberate
errors.
\item \texttt{r2\_least\_profile\_lp.py} (floating point; not a proof) runs
an LP that minimizes $P(n,n)$ under (S), (B), (C), (T) on a $(4n-1)^2$ grid.
The minimum agrees with $P^*(n,n)$ to $10^{-13}$ for $n\le12$.
\end{itemize}

\begin{thebibliography}{99}
\bibitem{OAI} OpenAI, \emph{An Upper Bound of $9/4$ for the Matrix
Multiplication Exponent}, OpenAI Math Release preprint, October 2, 2026,
\url{https://github.com/openai/math/blob/adc7f1241b42e322a6451854ab7e4b4c146bf78a/preprints/Matrix-Multiplication-Nine-Fourths-October-2-2026/build/paper.tex}.
\bibitem{EC} \texttt{eliottcassidy2000/math}, THM-4568, ``The growth lemma
behind omega $\le$ 9/4 (openai/math \#107) is exactly sharp'', commit
\texttt{450388a22f72}, 2026-10-07,
\url{https://github.com/eliottcassidy2000/math/blob/5966360a9c69b6bb0f3da39dc219f4a6c5dc3244/01-canon/theorems/THM-4568-the-nine-fourths-growth-lemma-is-exactly-sharp.md}.
\bibitem{SVA} \texttt{GasStationManager/SafeVerifyAgent},
\path{audits/mm94-intake/READ-F.md}, \S2d, commit \texttt{af3b9e1e8853},
2026-10-08.
\bibitem{SEL} S.~Navot et al., \texttt{selanavot/matrix-multiplication-all-fields},
$\omega\le9/4$ over arbitrary fields, commit \texttt{08481ef22bca}, 2026.
\bibitem{SUPP} This remark, \texttt{proofs/}: reports of the three
Claude--Codex rounds with the validity proofs, 2026-10-09.
\bibitem{Str88} V.~Strassen, The asymptotic spectrum of tensors,
\emph{J. Reine Angew. Math.} 384 (1988), 102--152.
\bibitem{CVZ} M.~Christandl, P.~Vrana, J.~Zuiddam, Universal points in the
asymptotic spectrum of tensors, \emph{J. Amer. Math. Soc.} 36 (2023), 31--79.
\bibitem{CVZb} M.~Christandl, P.~Vrana, J.~Zuiddam, Barriers for fast matrix
multiplication from irreversibility, \emph{Theory Comput.} 17 (2021).
\bibitem{AFLG} A.~Ambainis, Y.~Filmus, F.~Le Gall, Fast matrix
multiplication: limitations of the Coppersmith--Winograd method, STOC 2015.
\bibitem{Alm} J.~Alman, Limits on the universal method for matrix
multiplication, CCC 2019.
\bibitem{LG14} F.~Le Gall, Powers of tensors and fast matrix multiplication,
ISSAC 2014.
\bibitem{Wendel} J.~G.~Wendel, Note on the gamma function, \emph{Amer. Math.
Monthly} 55 (1948), 563--564.
\end{thebibliography}

\bigskip\noindent\emph{AI assistance.} Developed with AI assistance
(Anthropic Claude, Opus 5.5) under [AUTHOR NAME]'s direction, together with
an OpenAI Codex session (gpt-6-astra) for Section~\ref{sec:enlarged}.
Independent adversarial reviews were run by separate AI sessions and are not
human peer review.
\end{document}
